\section{Structural limits of two routes to a larger improvement}
\label{sec:structure}

The new consistency certificate improves a leading contraction coefficient.
Two tempting ways to improve the logarithmic order are to compress the
variable information further and to lower the initial tensor degree.
This section gives precise limits of particular implementations of those
ideas. None is a lower bound on the running time of 3-SAT.

\subsection{The variable-union pairing has full rank}

In the incidence encoding of the source reduction~\cite{openai}, the
variable-side tensor depends only on whether a variable appears in a selected
incidence monomial. If $S$ is its set of touched variables, summing that
monomial over Boolean assignments gives $2^{n-|S|}$. It is therefore natural
to ask whether a small algebra can retain just the information needed to
combine such sets.

\begin{theorem}[A fixed linear summary requires exponential dimension]
\label{thm:union-rank}
Let $\mathbb F$ be any field and put
\[
 A_n=\mathbb F[x_1,\ldots,x_n]/(x_i^2-x_i:1\le i\le n).
\]
For $S\subseteq[n]$, write $e_S=\prod_{i\in S}x_i$, and define the
linear functional $L:A_n\to\mathbb F$ by $L(e_S)=2^{n-|S|}$, with
the integer interpreted in $\mathbb F$.
Suppose $E:A_n\to V$ is linear, $\dim_{\mathbb F}V=d$, and a bilinear
form $b:V\times V\to\mathbb F$ satisfies
\[
 b(E(f),E(g))=L(fg)\qquad(f,g\in A_n).
\]
Then $d\ge2^n$.
\end{theorem}

\begin{proof}
The squarefree monomials $e_S$ form a basis and satisfy
$e_Se_T=e_{S\cup T}$. Consequently the matrix of the pairing
$(f,g)\mapsto L(fg)$, in this basis, is
\[
 H_{S,T}=2^{n-|S\cup T|}
 =\left(\begin{array}{cc}2&1\\1&1\end{array}\right)^{\otimes n}_{S,T}.
 \tag{U}
\]
Order subsets by their binary indicator strings. Each coordinate contributes
$2$ when both membership bits are zero, and $1$ otherwise, proving the
Kronecker identity. The one-coordinate matrix has the integer inverse
\[
 \left(\begin{array}{rr}1&-1\\-1&2\end{array}\right).
\]
Thus $H$ is invertible over every field, including characteristic two, and
has rank $2^n$. On the other hand, if $Q$ is the matrix of $E$ and $B$
the matrix of $b$, then $H=Q^{\mathsf T}BQ$ has rank at most $d$.
\end{proof}

This statement also applies to a linear encoding into a finite-dimensional
algebra whose multiplication is followed by a linear readout: their
composition is such a bilinear form. The basis inputs already arise from
3-CNFs of width at most three, since the conjunction of the positive unit
clauses indexed by $S$ has satisfaction polynomial $e_S$.

There is an equally direct decision version. For each
$\alpha\in\{0,1\}^n$, let $F_\alpha$ be the conjunction of unit clauses
fixing the assignment to $\alpha$. Then
\[
 \bigl(\mathbf1_{\SAT(F_\alpha\wedge F_\beta)}\bigr)_{\alpha,\beta}
 =I_{2^n}.
\]
Any fixed $d$-dimensional summaries whose conjunction decision is obtained
by a bilinear form over a field therefore require $d\ge2^n$, whether or not
the summary assignment is linear.

These are constraints on a fixed compositional representation. They do not
exclude adaptive restrictions, special reachable families, nonlinear
evaluation, or another summation method. In particular, they do not show
that the actual sparse parity-query family needs all $2^n$ dimensions.
They explain why merely identifying incidence subsets with variable unions
does not by itself supply the missing efficient algorithm.

\subsection{Budget-aware counting in the existing grouping reduction}

The source groups tests and variables separately into
$J=\lceil n^{1/4}\rceil$ groups~\cite{openai}. Increasing $J$ to nearly
$\sqrt n$ is already present in the corpus: the recursive refinement uses
$J=\lceil n^{0.495}\rceil$~\cite{recursive}, and the high-moment packet uses
$J=\lceil\sqrt n/\lceil\log_2n\rceil^2\rceil$~\cite{highmoment}.
The latter explicitly obtains initial degree
$\sqrt n\,\operatorname{polylog}(n)$. We credit that result; the refinement
below concerns its array-counting overhead.

Consider the source's sparse-query interface with degree at most three,
$a\le cn$ tests and at most $\lambda^c$ monomials per test, for fixed $c$,
where $\lambda=\lceil\log_2n\rceil$. After its high-incidence branching,
each test has at most $3\lambda^c$ incident variables, each free variable
has $O_c(\lambda^{c+1/2})$ incidences, and the number of branches is at most
$2^{n/\sqrt\lambda}$. The retained incidence-degree bound is
$K=3(a+2n)=O_c(n)$.

Partition both sides into $J$ groups balanced in cardinality. Let
$N_{gh}$ be the incidence count for a group pair, let
$N=\sum_{g,h}N_{gh}$, and let $B$ count the nonempty pairs.
Thus $B\le J^2$ and $N=O_c(n\lambda^c)$. The enumerated arrays satisfy
\[
 0\le d_{gh}\le N_{gh},\qquad \sum_{g,h}d_{gh}\le K.
\]

\begin{lemma}[Array and wire budgets]
\label{lem:array-budget}
The number $R$ of such arrays satisfies
\[
 R\le\min\left\{(K+1)^B,\ \binom{K+B}{B},\
                   \prod_{N_{gh}>0}(N_{gh}+1)\right\}.
 \tag{A}
\]
For each array the source's subset-index encoding has wire count $M$ and
initial maximum degree $\Delta_0$ satisfying
\begin{align*}
 M&\le B+\log_2\binom{N+K}{K},\qquad M\le N,\\
 \Delta_0&=O_c\bigl((n/J+1)\lambda^{c+1/2}\bigr).
\end{align*}
Its local-entry evaluation still takes
$\operatorname{poly}(n)2^{\lambda^{O_c(1)}}$ time and polynomial space.
\end{lemma}

\begin{proof}
The three bounds in (A) respectively discard the sum constraint, discard
the individual capacities and count weak compositions of a sum at most
$K$, and discard only the sum constraint while retaining the capacities.
For a fixed array, pair $(g,h)$ uses
$\lceil\log_2\binom{N_{gh}}{d_{gh}}\rceil$ wires. Their product of subset
counts is bounded by the number of monomials of total degree at most $K$
in $N$ indeterminates, namely $\binom{N+K}{K}$. Summing the rounding errors
gives the first wire bound. The inequality
$\lceil\log_2\binom{t}{j}\rceil\le t$ gives the second.

A group vertex has degree at most its incidence count; combine the balanced
group sizes with the test and variable incidence bounds. A test-group entry
is a product of independent per-test coefficients. Each is evaluated by the
source's finite-difference formula on one polylogarithmic scope.
The times for these evaluations add; they do not multiply. The
variable-group entry is the power of two recording untouched variables.
Subset decoding, ring arithmetic and all descriptions have polynomial size.
\end{proof}

The composition count improves the allowable polylogarithmic choice of $J$.
For $J\asymp\sqrt n/\lambda^A$ with fixed $A>0$, it gives
\[
 \log_2R
 \le B\log_2\!\left(\frac{\mathrm e(K+B)}B\right)
 =O_c\!\left(\frac{n\log\lambda}{\lambda^{2A}}\right).
\]
Hence $A>1/4$ suffices for $o(n/\sqrt\lambda)$ overhead; the crude
$(K+1)^{J^2}$ estimate would require $A>3/4$.

For an exact integer choice, put $\nu=\lceil\log_2\lambda\rceil$ and take
\[
 J=\max\{j\in\mathbb Z_{\ge0}:
               j^2\lceil\sqrt\lambda\rceil\nu^2\le n\}.
\]
Treat the finite range with $J<1$ directly. Then
\[
 J=\Theta\!\left(\frac{\sqrt n}{\lambda^{1/4}\nu}\right),\quad
 \log_2R=O_c\!\left(\frac n{\sqrt\lambda\,\nu}\right),\quad
 \Delta_0=O_c(\sqrt n\,\lambda^{c+3/4}\nu).
\]
The wire bound remains $O_c(n\nu)$, and $B=o(n)$.
Arrays can be generated by a depth-first search with a remaining-degree
counter, using polynomial work per array and polynomial space.
The extra exponent outside contraction is
$O(n/\sqrt\lambda)+O_c(n/(\sqrt\lambda\,\nu))$, including the original
high-incidence branching.

Thus the improvement is in polylogarithmic slack, while the initial degree
exponent remains $1/2$. A schedule accepting
$s^{1/2}(\log s)^{O(1)}$ degrees may use that bound directly. A theorem
requiring exactly $\Delta_0\le s^\gamma$ may instead use any fixed
$\gamma>1/2$, unless its parameter construction explicitly absorbs the
polylogarithmic factor. The high-moment packet supplies such an integer
budget~\cite{highmoment}. This reasoning does not establish the stronger
degree promise $s^{1/12}$ for arbitrary-density queries.

\subsection{Small groups can force exponentially many degree arrays}

The next result strengthens the warning against simply choosing many more
groups. Its quantifiers allow the partition to be chosen after seeing the
graph.

\begin{theorem}[A graph-adaptive grouping obstruction]
\label{thm:group-obstruction}
Fix $0<\delta<1/2$ and an integer $d$ with
$d(1-2\delta)>3$. For every sufficiently large $n$ there exists a bipartite
graph with $n$ vertices on each side and maximum degree at most $d$ such
that every partition of each side into groups of size at most $n^\delta$
has more than $n$ nonempty group pairs.
\end{theorem}

\begin{proof}
An ordered $d$-tuple of permutations of $[n]$ specifies $d$ perfect
matchings. Forget the colors and repeated edges; the resulting simple
graph has maximum degree at most $d$. There are
\[
 (n!)^d=\exp(dn\ln n+O_d(n))
\]
tuples. We upper-bound the tuples whose graph has a prohibited partition.

Label groups by numbers in $[n]$. There are at most $n^{2n}$ pairs of
labeled partitions. A support of at most $n$ group pairs can be chosen in
at most
\[
 (n+1)\binom{n^2}{n}=\exp(n\ln n+O(n))
\]
ways. If groups have size at most $q=\lfloor n^\delta\rfloor$, this support
allows at most $Q=nq^2\le n^{1+2\delta}$ individual vertex pairs.

The projected graph has at most $dn$ distinct edges. Its edge set has at
most
\[
 \sum_{t=0}^{dn}\binom Qt
 \le\binom{Q+dn}{dn}
 =\exp(2d\delta n\ln n+O_d(n))
\]
possibilities. The inequality follows by representing a subset of size at
most $dn$ as a multiset of size exactly $dn$ after adjoining a dummy
element. For a fixed graph of maximum degree at most $d$, at most
$d^{dn}$ permutation tuples project to it: at each left vertex each color
chooses one of at most $d$ neighbors. The permutation requirement only
reduces that count.

Multiplying these bounds gives
$\exp((3+2d\delta)n\ln n+O_d(n))$ bad tuples. This is less than
$(n!)^d$ for large $n$ by the assumed strict inequality. Some tuple
therefore gives a graph with the stated property.
\end{proof}

\begin{corollary}[The total-degree constraint does not remove this cost]
\label{cor:array-obstruction}
For the graphs in Theorem~\ref{thm:group-obstruction}, every permitted
grouping has at least $2^n$ valid degree arrays whenever $K\ge n$.
\end{corollary}
\begin{proof}
Choose any $n$ nonempty group pairs, assign $d_{gh}\in\{0,1\}$
independently on them, and assign zero elsewhere. All capacities are
respected and the total is at most $n$.
\end{proof}

For example, $\delta=1/10$ and $d=4$ satisfy the strict inequality.
Group sizes $n^{1/12}\operatorname{polylog}(n)$ are eventually at most
$n^{1/10}$. The result therefore blocks a direct attempt to obtain the
$s^{1/12}$ promise by uniformly shrinking all groups while continuing to
enumerate the same arrays.

The graphs fit the generic sparse-query interface: take one polynomial
per left vertex equal, over $\mathbb F_2$, to the sum of its neighboring
variables. These have degree one and scopes of size at most $d$, so there
is no high-incidence branching and $K\ge n$. This example is easy by
Gaussian elimination. Its role is precisely to show that the literal
outer enumeration can be expensive even on easy inputs unless further
algebraic structure is used.

Finally, small group cardinality is a \emph{sufficient} way of guaranteeing
small encoded degree, not a necessary one. A specific array with
$d_{gh}=0$ or $N_{gh}$ can have zero wires on that pair even when the group
is large. The theorem neither rules out exploiting these cases nor
lower-bounds the number of arrays with nonzero coefficient. It rules out
only the stated combination of uniformly small groups and exhaustive
enumeration of combinatorially allowed arrays.

The accompanying \texttt{check\_structures.py} verifies the integer
Kronecker inverse, its reductions in characteristic two and odd
characteristics, direct small Boolean counts, and capacity-constrained
array counts against independent exhaustive enumeration. Deliberately
corrupted matrices and counts must be rejected. These checks support the
finite identities; the quantified graph obstruction follows from the
proof above.

